\section{Prospective supplement: the source-expansion and transport-bridge programme}
\label{sec:prospective-source-expansion}

Nothing in this section is result-bearing. It records an outcome-free census and
the state of a transport bridge that has not closed: the bridge stands at 76/80
and the programme has adjudicated \textbf{zero} natural pairs. It is reported
because an unrun arm and an absent section look identical to a reader, and only
one of them is a disclosure.

It sits outside the results deliberately. The counts below are lower bounds on
observable candidates, not adjudicated pairs, and no claim in this paper rests
on them. They describe what a future naturalistic comparison would have to draw
from, and what still blocks it.

\subsection{Outcome-free source-expansion census}
The naturalistic successor was not executed.  We instead ran a separate,
outcome-free Zenodo feasibility census for its article--data (M5) and
article--software (M6) source routes.  The frozen V1 request failed before any
census because Zenodo rejected its unauthenticated page size of 200 with HTTP
400; that transport terminal is retained.  A distinct V2 successor disclosed
and excluded every record identity seen in the transport/schema probes, used
nine pages of 25 records for each of eight domain--mechanism queries, required
an accepted open record licence and public file-link evidence, and counted only
accepted relations whose target was publication-typed.  All 72 responses passed
the frozen schema and pagination checks, with no within-query duplicates.

V2 retained 173 query-specific candidates representing 98 unique Zenodo
records.  The publication-linked counts for the four M5 data queries were 37,
44, 29 and 30; the corresponding M6 software counts were 6, 10, 10 and 7.
Thus none of the eight cells reached the unchanged pre-adjudication signal gate
of 48.  The contrast diagnoses different source bottlenecks: the M5 route is
near but below quota, whereas M6 has many accepted typed relations but few
publication-typed targets.  The terminal is a source-cell shortfall requiring a
disjoint provider route, not a natural-pair, scientific-performance or model
outcome.

\begin{table}[t]
\centering
\small
\caption{Frozen Zenodo V2 publication-typed source signals after licence,
public-file and pre-freeze-exclusion gates.  The threshold is 48 in every cell;
none passes.  Counts are discovery candidates, not adjudicated pairs.}
\label{tab:p4-zenodo-source-signals}
\begin{tabular}{lrr}
\toprule
Domain & M5 article--data & M6 article--software \\
\midrule
Earth/environment & 37 & 6 \\
Life/biomedical & 44 & 10 \\
Physical/engineering & 29 & 10 \\
Scientific software & 30 & 7 \\
\bottomrule
\end{tabular}
\end{table}

We then froze a provider-disjoint DataCite M5 successor without changing the
four domains, 48-signal threshold, exact CC BY 4.0/CC0 declaration,
HTTPS-content-URL requirement, publication-typed relation, or rights boundary.
The first DataCite identity returned 4,000 valid metadata rows but omitted the
client relationship required for provider disjointness; all rows therefore
failed closed before rights or relation screening. Its terminal
\idt{P4_DATACITE_M5_CENSUS_V1_CANNOT_CHECK_TRANSPORT_OR_SCHEMA} is retained.

A distinct V2 mechanics repair bound and excluded all 3,479 unique DOI
identities disclosed by V1, requested the client relationship, retained Zenodo
source-family exclusions, and corrected only a literal documentation assertion.
All four response schemas, pages and evidence assertions passed. Of 4,000 raw
rows, 968 were excluded by the frozen V1 disclosure, leaving 3,032
provider-disjoint rows; 2,719 carried an exact CC BY 4.0 or CC0 declaration.
None exposed an HTTPS \texttt{contentUrl}, so no row reached publication-relation
adjudication and no candidate was admitted. The combined M5 counts therefore
remain 37, 44, 29 and 30, while the separately preserved M6 counts remain 6,
10, 10 and 7. The terminal is
\idt{P4_NATURAL_PAIR_SOURCE_CELL_SHORTFALL__EXPAND_DISJOINT_SOURCE_UNIVERSE}.
The zero relation count is a sequential not-reached count, not evidence that
DataCite records globally lack publication relations. V2 repairs transport
mechanics only; it is not a positive source, case or model result.

A separately frozen multi-provider expansion then audited Dryad, Figshare,
Harvard Dataverse and DataCite for both M5 and M6. It retained 1,263 successful
raw API responses, 26 provider-provenance captures and 2,371 candidate rows,
while counting at most one candidate per concept DOI and linked publication
DOI. Mirrors and metadata aggregators were not treated as independent
providers. Admission required an exact accepted item licence, a structured
publication relation and an unrestricted direct-download identity.

Provider transport was incomplete. Only 659 of 2,463 frozen Figshare full
records and 512 of 1,540 Harvard Dataverse full records were captured; 1,804
and 1,028 identities, respectively, remain frozen rather than replaced. Dryad
returned 1,200 metadata records, but anonymous download probes returned HTTP
401 and exact version/file lineage remained undetermined. DataCite supplied
2,958 discovery DOIs but zero counted units because the repository-specific
record and file authority was not bound. The strict deduplicated counts in
Table~\ref{tab:p4-multiprovider-source-lower-bounds} are therefore lower-bound
diagnostics, not confirmed cell deficits, and the eight-cell gate is not
evaluable.

\begin{table}[t]
\centering
\small
\caption{Strict deduplicated public-source candidates observed before transport
completion. Counts are lower bounds, not adjudicated natural pairs or confirmed
quota failures.}
\label{tab:p4-multiprovider-source-lower-bounds}
\begin{tabular}{lrr}
\toprule
Domain & M5 article--data & M6 article--software \\
\midrule
Earth/environment & 6 & 2 \\
Life/biomedical & 10 & 0 \\
Physical/engineering & 4 & 1 \\
Scientific software & 3 & 6 \\
\bottomrule
\end{tabular}
\end{table}

The source-feasibility terminal is
\idt{P4_NATURAL_PAIR_SOURCE_TRANSPORT_CANNOT_CHECK}. Natural-pair identity,
author-lineage disjointness, same-claim preservation and scientific
resolvability still require outcome-blind external adjudication. No case,
model-performance or superiority result follows.

We next opened a separately frozen M6 provider route without altering any V3
identity. Its first Crossref page fixed 200 unique Journal of Open Source
Software publications. All 200 JOSS pages transported; 191 exposed a labelled
GitHub repository relation, 180 repositories exposed a latest release whose tag
resolved to an immutable commit, 118 bound an accepted licence and licence-blob
SHA at that exact tag, and 120 received one domain under the unchanged
outcome-free token rule. The intersection contains 80 unique publication-DOI
and repository concepts: 5 Earth/environment, 7 life/biomedical, 62 scientific
software and 6 physical/engineering.

Those 80 are provider-qualified concepts for a later version-link and
natural-pair audit, not eligible units. The JOSS page relates the paper to a
repository, but this lane did not establish that the current GitHub release tag
is the version evaluated by the paper. It therefore established 0 exact
paper-to-release relations, 0 author-lineage adjudications and 0 natural pairs.
The bounded V3 packet also omits its row identities, so cross-provider
publication, repository and artifact deduplication is \textsc{CannotCheck}.

Even before those unresolved gates, the no-overlap union is only an optimistic
diagnostic. Earth, life and physical each reach 7/48. Scientific software
reaches 68/48, but JOSS/GitHub is one provider family and the disjoint Figshare
side contains only 6/8 units. Thus neither surplus repositories nor releases,
tags, assets, files, commits, versions, search hits or API responses can satisfy
the frozen allocation gate. The V4 terminal is
\idt{P4_M6_JOSS_GITHUB_BOUNDED_TRANSPORT_COMPLETE__EXACT_PUBLICATION_RELEASE_VERSION_RELATION_AND_AUTHOR_LINEAGE_CANNOT_CHECK__M6_CELL_FRAME_NOT_READY}.
The programme terminal remains exactly
\idt{P4_NATURAL_PAIR_SOURCE_TRANSPORT_CANNOT_CHECK}; no quota, source-frame,
naturalistic or performance claim is promoted.

The no-extension V5 bridge then retained exactly those 200 JOSS publication
DOIs and followed each labelled Software archive DOI through DataCite and
source-native archive metadata. Of 200 archives, 103 exposed an explicit tag or
commit identity for the same JOSS repository and 101 resolved to an immutable
Git commit. After also requiring the original V4 provider-qualified identity,
one distinct archive-version/concept relation, accepted exact archive and
commit rights, unchanged repository/domain identity, source-native archive
transport, and exact deduplication against 9/9 hash-recovered V3 M6 identities,
39/80 pass the publication--archive-version--tag--commit bridge: 3 Earth, 4
life, 30 scientific-software and 2 physical concepts. Each is one unique
publication-DOI/repository concept; files, tags, releases and commits do not
multiply the scientific unit.

This closes an exact public transport relation for those 39, not natural-pair
eligibility. The deduplicated V3-strict plus V5-exact candidate totals are
5/48, 4/48, 36/48 and 3/48; the Figshare replication sides are 2/8, 0/8, 6/8
and 1/8. All four gates fail in this exact transported frame; this does not
convert the incomplete V3 provider census into a global population deficit.
Author-lineage adjudications and
natural pairs remain 0. Of the 80 V4-qualified concepts, 32 lack an
archive-explicit same-repository tag/commit relation, 12 lack a distinct
archive-version/concept relation and 3 lack accepted exact archive rights;
overlap among failures leaves 41 incomplete. The V5 terminal is
\idt{P4_M6_JOSS_ARCHIVE_EXACT_VERSION_BRIDGE_PARTIAL_39_OF_80__M6_CELL_FRAME_AUTHOR_LINEAGE_AND_NATURAL_PAIR_CANNOT_CHECK}.
It preserves the V4 and programme terminals and supplies no case, system
outcome, naturalistic comparison or superiority evidence.

A prospectively separated V6 repair retained exactly the same 41 failed
identities and added or replaced zero publication DOIs. Exact normalized
source-archive/GitHub immutable-commit manifests repaired 12 identities; five
additional identities were repaired by equality between a qualified
source-native Software Heritage directory and the immutable Git commit root
tree. Exact version/concept identity and accepted software rights at both the
archive and commit remained mandatory. V6 repairs 17/41 and raises the exact
bridge from 39/80 to 56/80: 4 Earth, 4 life, 45 scientific-software and 3
physical concepts. Files, versions, tags, commits, requests and manifests add
zero scientific units.

The 24 unresolved identities comprise 12 exact-publication-version failures,
2 exact-archive-rights failures and 10 archive-to-commit content-identity
failures under a mutually exclusive primary-cause assignment. Deduplicated
V3-strict plus V6-final totals are 6/48, 4/48, 51/48 and 4/48; the unchanged
Figshare replication sides are 2/8, 0/8, 6/8 and 1/8. Scientific software
clears its total and JOSS-primary quotas but not its source-disjoint replication
quota, while the other three domains remain below both gates. Thus every full
cell fails, author-lineage adjudications and natural pairs remain 0, and the
V6 terminal is
\idt{P4_M6_JOSS_SAME_IDENTITY_CONTENT_ADDRESSED_BRIDGE_REPAIR_PARTIAL_17_OF_41__FINAL_EXACT_56_OF_80__AUTHOR_LINEAGE_AND_NATURAL_PAIR_CANNOT_CHECK}.
This is development transport evidence with a disclosed pre-freeze route probe,
not confirmatory, global-provider, natural-pair, performance or superiority
evidence. The programme terminal remains unchanged.

A targeted V7 successor then revisits only those 24 unresolved publication
identities. It adds no publication DOI and does not count versions, tags,
commits, files or requests as new units. Requiring an exact JOSS review relation,
a unique publication-time archive child, a source-native tag resolved to a
40-hex commit, archive-authenticated content/origin identity, provider-bound
checksums, and accepted rights at archive and commit repairs 14/24. The exact
bridge rises from 56/80 to 70/80.

Ten identities remain unresolved under a mutually exclusive primary partition:
three archive-to-commit authenticated-identity failures, two archive-software-
rights failures, two exact archive-version DOI failures, and three exact
tag-to-immutable-commit failures. Editorial assertions do not close these gates,
and the packet performs zero replacement-publication, author-lineage or
natural-pair adjudications. Its programme terminal is therefore
\idt{P4_NATURAL_PAIR_SOURCE_TRANSPORT_CANNOT_CHECK}. V7 is targeted development
transport evidence, not a naturalistic case, model outcome or superiority
result.

A still narrower V8 pass targets only those ten disclosed identities. Ten of
ten provider archives are fetched once and checksum-verified; no publication
is replaced and no file, version or request is counted as a scientific unit.
Provider-native version, archive provenance and licence evidence repairs PSD,
FAIRLinked and SEPARATE, raising the exact bridge from 70/80 to 73/80 in
36.011 measured seconds. Seven identities remain fail-closed, each with one
exact missing authoritative edge. The pass adds no author-lineage
adjudication, natural pair, replication unit, comparator execution or outcome,
so the programme terminal remains
\idt{P4_NATURAL_PAIR_SOURCE_TRANSPORT_CANNOT_CHECK}.

A V9 same-identity successor subsequently repairs two of those seven edges by
exact normalized checksum-bound archive--immutable-revision manifest equality,
raising the bridge from 73/80 to 75/80.  Of the five residual identities, a
frozen V10 provider-authority pass leaves four unchanged and its target-specific
test for the fifth fails as registered.  A separately named, explicitly
outcome-informed V10B protocol is then frozen before new archive download or
comparison.  For that fifth identity, an immutable later revision is bound by
both GitHub and the repository's Software Heritage origin snapshot to the same
Git tree as the publication archive; its checksum-bound Zenodo ZIP and GitHub
codeload ZIP normalize to the same 165 paths, with zero missing, extra or
differing files and byte-identical MIT licences.  This exact-content witness
raises the bridge to 76/80 without identifying the later revision as the
deleted publication-version commit.  Four authoritative edges remain.  V9,
V10 and V10B add no publication unit, lineage adjudication, natural pair,
replication unit, comparator execution or outcome, so the programme terminal
remains \idt{P4_NATURAL_PAIR_SOURCE_TRANSPORT_CANNOT_CHECK}.

V11 revisits exactly the four remaining edges and closes none.  The adverse
result separates provider-state changes from repeated querying: index 36 needs
a corrected checksum-bound archive whose root is the named publication commit;
indices 133 and 185 need exact-file signed PyPI provenance to their full
commits; and index 91 needs either exact payload equality or provider-native
signed build provenance.  Its cumulative bridge therefore remains 76/80.

V12 tests the index-91 causal split.  The GitHub-authenticated immutable tree
tracks \texttt{java/jPAI.jar} with SHA-256
\idt{ff49482838e3a761913df327a48e819d4f9e552bfeab29655a82675c60c47162}.
A frozen read-only projection yields exactly 106 class members, and their
complete \texttt{java/bin/} multiset equals the 106 archive-only class files
with zero missing, extra or differing bytes; all have Java class major version
61.  This is a positive exact packaged-byte result, not a reproduced build.
The pinned Java~17 OCI compiler did not execute because no Docker daemon was
available and the disposable Colima fallback ended with the exact terminal
\idt{no_space_left_on_device}; no substitute compiler was used.  Moreover, the
validly signed revision tree records both \texttt{cpp/PAIpp.exe} and
\texttt{run.sh} as mode 100644, while the archive records 0755, and no signed
release or artifact attestation authorizes either change.  V12 therefore closes
zero edges and leaves the bridge at 76/80.

